\subsection{平面旋转的补充}

本号中，V 表示一个二维实向量空间，它配备了一个标量积（即非退化、正定的对称双线性型）和一个定向。角的度量取底 \(2\pi\)；半直线（分别为直线）之间的角的主度量

因而 \(\theta\) 是满足 \(0 \leqslant \theta < 2\pi\)（分别为 \(0 \leqslant \theta < \pi\)）的实数（《一般拓扑》第 VIII 章 §2 第 3 号和第 6 号）。滥用术语，对于任意实数 \(\theta\)，我们用 \(\theta\) 表示度量为 \(\theta\) 的角，并记 \(\rho_{\theta}\) 为角度为 \(\theta\) 的旋转（《代数》第 IX 章 §10 第 3 号）。

命题 6. 设 s 是 V 中关于直线 D 的正交反射。若 \(\Delta\) 和 \(\Delta'\) 是从原点出发的两条半直线（分别为经过原点的两条直线），则有

\[
(s (\widehat {\Delta}), s (\Delta^ {\prime})) \equiv - (\widehat {\Delta}, \Delta^ {\prime}) \pmod {2 \pi} \text {(resp. (mod. } \pi \text {))}.
\]

设 \(u\) 是将 \(\Delta\) 变换为 \(\Delta'\) 的旋转。由于 \(su\) 是 \(V\) 的行列式为 \(-1\) 的正交变换，它是反射，因而 \((su)^2 = 1\)。所以，\(u^{-1} = sus^{-1}\) 将 \(s(\Delta)\) 变换为 \(s(\Delta')\)，从而命题得证。

推论. 设 D 和 \(D'\) 是 V 中的两条直线，令 \(\theta\) 为角 \((\widehat{\mathrm{D},\mathrm{D}'}\) ) 的一个度量。则 \(s_{D'}s_{D} = \rho_{2\theta}\)。

我们知道 \(s_{D'} s_{D}\) 是一个旋转，因为它的行列式为 1。设 \(\Delta\) 和 \(\Delta'\) 是由 D 和 \(D'\) 承载的从原点出发的两条半直线。我们有

\[
\begin{array}{r l} (\widehat {\Delta , s _ {\mathrm {D^ {\prime}}} s _ {\mathrm{D}}} (\Delta)) & \equiv (\widehat {\Delta , s _ {\mathrm {D^ {\prime}}}} (\Delta)) \equiv (\widehat {\Delta , \Delta^ {\prime}}) + (\Delta^ {\prime}, \widehat {s _ {\mathrm {D^ {\prime}}}} (\Delta)) \\ & \equiv (\widehat {\Delta , \Delta^ {\prime}}) + (s _ {\mathrm {D^ {\prime}}} (\widehat {\Delta^ {\prime}}), s _ {\mathrm {D^ {\prime}}} (\Delta)) \\ & \equiv (\widehat {\Delta , \Delta^ {\prime}}) - (\widehat {\Delta^ {\prime} , \Delta}) \equiv 2 (\widehat {\Delta , \Delta^ {\prime}}) \pmod {2 \pi}, \end{array}
\]

因此推论得证。

现在设 \(\Delta\) 和 \(\Delta'\) 是 \(V\) 中满足下列条件的两条半直线：

\[
\Delta \neq \Delta^ {\prime} \text { and } \Delta \neq - \Delta^ {\prime};
\]

并令 \(s\) 和 \(s'\) 为关于包含 \(\Delta\) 和 \(\Delta'\) 的直线 D 和 D' 的正交反射。令 \(\theta\) 为角 \((\widehat{\mathrm{D},\mathrm{D}'})\) 的主度量。若 \(\theta \in \pi\mathbf{Q}\)，记 \(m\) 为满足 \(\geqslant 1\) 且 \(m\theta \in \pi\mathbf{Z}\) 的最小整数。若 \(\theta \notin \pi\mathbf{Q}\)，置 \(m = \infty\)。令 W 为由 \(s\) 和 \(s'\) 生成的群。

命题 7. 群 W 是二面体群（《代数》第 IV 章 §1 第 2 号），阶为 \(2m\)。它由旋转 \(\rho_{2n\theta}\) 和乘积 \(\rho_{2n\theta}s\) 构成，其中 \(n \in \mathbf{Z}\)。W 的元素对 D 和 D' 的变换，正是将 D 通过旋转 \(\rho_{n\theta}\) 所作的变换，其中 \(n \in \mathbf{Z}\)。

命题 6 的推论表明 \(s's\) 的阶为 \(m\)，这给出第一个断言。于是，\(W\) 的元素形如 \((s's)^n = \rho_{2n\theta}\) 或 \((s's)^n s = \rho_{2n\theta}s\)；最后一个断言由此得出，因为 \(D' = \rho_\theta(D)\)。

推论. 令 C 为由从原点出发的开半直线 \(\Delta_{1}\)（满足 \(0 < (\widehat{\Delta, \Delta_{1}}) < \theta\)）的并所成的开角扇区。则 D 或 \(\mathrm{D}'\) 经 W 的元素变换所得的集合均不与 C 相交，当且仅当 \(m\) 有限且

\[
\theta = \pi / m.
\]

若 \(m = \infty\)，则集合 \(n\theta\)（\(n \in Z\)）在 R/2πZ 中的像稠密（《一般拓扑》第 VII 章 §1，命题 11）；因而，D 由 W 中元素变换所得的并在 V 中稠密并与 C 相交。若 m 有限且 \(\theta = k\pi/m\)，其中 1 \textless{} k \textless{} m，且整数 k 与 m 互质，则存在整数 h 使 \(hk \equiv 1 \mod m\)；于是 \((\widehat{\mathrm{D}, \rho_{h\theta}(\mathrm{D}))} \equiv \pi/m \pmod{\pi}\)，且 \(\rho_{h\theta}(\mathrm{D})\) 与 C 相交。这说明该条件是必要的。反之立即可见。

注. 假设 \(m\) 有限且 \(\theta = \pi / m\)。若 \(n \in \mathbf{Z}\)，令 \(C_n\) 为从 0 出发且满足下列条件的开半直线 \(\Delta_1\) 的并：

\[
n \theta <   (\widehat {\Delta , \Delta_ {1}}) <   (n + 1) \theta .
\]

\(C_{n}\)（其中 \(-m \leqslant n < m\)）是连通开子集，它们构成 \(E - \bigcup_{n} D_{n}\) 的一个划分（其中 \(D_{n} = \rho_{n\theta}(D)\)）。因此，这些就是由 m 条直线 \(D_{n}\)（\(1 \leqslant n \leqslant m\)）在 E 中确定的室。我们有 \(C_{2k} = \rho_{2k\theta}(C)\) 且 \(C_{2k-1} = \rho_{2k\theta}s(C)\)。此外，\(C_{n} = C\) 当且仅当 \(n \in 2mZ\)。因此，群 W 简单传递地置换室 \(C_{n}\)。

最后我们证明：若 \(w \in W\)，使得室 C 与 \(w(\mathbf{C})\) 位于直线 D 的相反两侧，则 \(l(sw) = l(w) - 1\)（长度是相对于 \(S = \{s, s'\}\) 取的）。事实上，假设蕴含 \(w(\mathbf{C}) = \mathbf{C}_n\)，其中 \(-m \leqslant n < 0\)。若 n = -2k，则 \(w = (ss')^k\) 且 \(sw = s'(ss')^{k-1}\)，所以 \(l(w) = 2k\) 且 \(l(sw) = 2k - 1\)（《代数》第 IV 章 §1 第 2 号，注）。若 \(n = -2k + 1\)，则

\[
w = \left(s s ^ {\prime}\right) ^ {k - 1} s \quad \text { and } \quad s w = \left(s ^ {\prime} s\right) ^ {k - 1},
\]

从而 \(l(w) = 2k - 1\) 且 \(l(sw) = 2k - 2\)。证毕。

